In this note, we report on a study using the CERN ProtoDUNEs as an opportunistic beam-dump experiment with respect to the upstream SPS T2 secondary beam target. This set-up can be used to search for and to study light particles expected to be produced in many Beyond the Standard Model (BSM) extensions. In a first stage we study the detection of beam-related neutrino interactions in the detectors. The study is based on a small subset of ProtoDUNE data collected in a parasitic mode. 

The ProtoDUNE Horizontal Drift (PD-HD) and ProtoDUNE Vertical Drift (PD-VD) detectors, located in the North Area at CERN, are two prototype Liquid Argon Time Projection Chambers (LArTPCs) constructed to validate the technologies proposed for the DUNE far detectors. During North Area operations, the T2 target is struck by 400 GeV protons from the SPS, generating a highly energetic and intense flux of mesons and other secondary particles within the target region and the dump material. Very recently, it was pointed out \cite{coloma2024new} that the ProtoDUNE detectors are coincidentally located along the beam direction of T2 and may be able to search for new physics signals from weakly interacting particles produced in the target and in the dump material. The proposed search would make use of existing infrastructure and may be carried out with minimal operational costs. However, it faces several challenges. Most notably, the ProtoDUNE modules operate on surface and are consequently exposed to an intense flux of cosmic rays, which requires a dedicated trigger. In addition, standard neutrinos are also produced in the T2 target area from the decay of unstable mesons, constituting a relevant background which needs to be well understood and characterized a priori. The analysis presented in this technical note constitutes an initial step in this direction and serves as a proof of concept. Before any search for new BSM can be attempted, it is imperative to demonstrate that standard weakly interacting
particles (namely neutrinos) coming from the T2 target can be detected, which constitutes the focus of this note.

The ProtoDUNE Horizontal Drift cryostat contains approximately 700 tons of liquid argon, and its TPC features the full 3.6~m drift distance corresponding to the DUNE Far Detector Module 1 (FD1-HD) design. At the nominal operating drift field of 500 V/cm, this corresponds to an electron drift velocity of approximately 1.6 mm/$\mu$s, giving a maximum drift time of approximately 2.3 ms across the full 3.6 m drift volume. The construction was completed in the second quarter of 2018, with data taking beginning shortly thereafter. The detector was upgraded later with FD pre-production components and equipped with the final FD data acquisition (DAQ) system. The detector was successfully commissioned, and approximately ten weeks of beam data were collected during the summer of 2024. In this technical note, we report the first observation of potentially accelerator-produced neutrino interaction candidates in PD-HD, from additional data recorded parasitically at the conclusion of the 2024 ProtoDUNE Horizontal Drift campaign using a dedicated trigger and offline software filter. Our results exhibit a clear excess above background along the SPS/T2 beam direction, in accord with simulation predictions.

These studies form a first step in searches for new physics, and also yield a sample of neutrino-candidate events detected in the LArTPCs (with energies roughly above 10 GeV), that can be used for further studies. Looking ahead, explicit BSM search algorithms can be deployed, and such studies will be instrumental to continue demonstrating the LArTPC technology's capabilities as an interesting option for future BSM searches using high-intensity beam-dump facilities~\cite{Ferrillo:2023hhg}.